\documentclass[11pt,a4paper]{article}
\usepackage{fontspec}
\defaultfontfeatures{Ligatures=TeX}
\setmainfont{Latin Modern Roman}
\usepackage{unicode-math}
\setmathfont{Latin Modern Math}
\usepackage{amsmath,booktabs,graphicx,geometry,natbib,hyperref}
\usepackage[protrusion=false]{microtype}
\usepackage{xcolor,caption,fancyhdr,setspace,placeins,array}
\newcolumntype{P}[1]{>{\raggedright\arraybackslash}p{#1}}
\geometry{left=26mm,right=26mm,top=24mm,bottom=25mm,headheight=14pt,headsep=9mm}
\definecolor{ink}{HTML}{213F50}
\hypersetup{colorlinks=true,linkcolor=ink,citecolor=ink,urlcolor=ink,pdftitle={Geopolitical Risk and Returns in Rare Earth and Strategic Metal Markets},pdfauthor={}}
\captionsetup{font=small,labelfont=bf,justification=justified,singlelinecheck=false,skip=7pt}
\setstretch{1.03}
\setlength{\parindent}{1.15em}
\setlength{\parskip}{1.5pt}
\setlength{\textfloatsep}{15pt plus 3pt minus 3pt}
\setlength{\floatsep}{12pt plus 2pt minus 2pt}
\setlength{\intextsep}{13pt plus 2pt minus 2pt}
\setlength{\bibsep}{4pt}
\setcounter{topnumber}{2}
\setcounter{bottomnumber}{2}
\renewcommand{\topfraction}{.88}
\renewcommand{\bottomfraction}{.75}
\renewcommand{\textfraction}{.10}
\renewcommand{\floatpagefraction}{.77}
\pagestyle{fancy}
\fancyhf{}
\fancyhead[L]{\small\itshape Geopolitical Risk and Strategic Metals}
\fancyhead[R]{\small 2010–2026}
\fancyfoot[C]{\small\thepage}
\renewcommand{\headrulewidth}{.3pt}
\raggedbottom
\emergencystretch=1.3em
\newcommand{\OosRemxLowCagr}{6.87}
\newcommand{\OosRemxLowDrawdown}{12.16}
\newcommand{\OosRemxLowFinal}{17,739}
\newcommand{\OosRemxLowTime}{4.0}
\newcommand{\OosRemxLowAlpha}{4.01}
\newcommand{\OosRemxLowAlphaP}{0.176}
\newcommand{\OosRemxHighCagr}{-7.89}
\newcommand{\OosRemxHighDrawdown}{73.34}
\newcommand{\OosRemxHighFinal}{4,920}
\newcommand{\OosRemxHighTime}{95.9}
\newcommand{\OosRemxHighAlpha}{-12.81}
\newcommand{\OosRemxHighAlphaP}{0.213}
\newcommand{\OosRemxAvoidCagr}{11.71}
\newcommand{\OosRemxAvoidDrawdown}{56.95}
\newcommand{\OosRemxAvoidFinal}{26,002}
\newcommand{\OosRemxAvoidTime}{55.1}
\newcommand{\OosRemxAvoidAlpha}{8.55}
\newcommand{\OosRemxAvoidAlphaP}{0.358}
\newcommand{\OosRemxMonthlyMACagr}{10.99}
\newcommand{\OosRemxMonthlyMADrawdown}{41.03}
\newcommand{\OosRemxMonthlyMAFinal}{24,594}
\newcommand{\OosRemxMonthlyMATime}{35.7}
\newcommand{\OosRemxMonthlyMAAlpha}{7.92}
\newcommand{\OosRemxMonthlyMAAlphaP}{0.349}
\newcommand{\OosRemxDailyMACagr}{11.27}
\newcommand{\OosRemxDailyMADrawdown}{52.01}
\newcommand{\OosRemxDailyMAFinal}{25,124}
\newcommand{\OosRemxDailyMATime}{38.8}
\newcommand{\OosRemxDailyMAAlpha}{8.41}
\newcommand{\OosRemxDailyMAAlphaP}{0.342}
\newcommand{\OosRemxMomentumCagr}{7.85}
\newcommand{\OosRemxMomentumDrawdown}{49.16}
\newcommand{\OosRemxMomentumFinal}{19,185}
\newcommand{\OosRemxMomentumTime}{42.6}
\newcommand{\OosRemxMomentumAlpha}{6.74}
\newcommand{\OosRemxMomentumAlphaP}{0.456}
\newcommand{\OosRemxMarketCagr}{1.68}
\newcommand{\OosRemxMarketDrawdown}{68.35}
\newcommand{\OosRemxMarketFinal}{11,543}
\newcommand{\OosRemxMarketTime}{44.6}
\newcommand{\OosRemxMarketAlpha}{-1.82}
\newcommand{\OosRemxMarketAlphaP}{0.838}
\newcommand{\OosRemxGlobalForecastCagr}{16.15}
\newcommand{\OosRemxGlobalForecastDrawdown}{40.25}
\newcommand{\OosRemxGlobalForecastFinal}{36,390}
\newcommand{\OosRemxGlobalForecastTime}{42.5}
\newcommand{\OosRemxGlobalForecastAlpha}{13.60}
\newcommand{\OosRemxGlobalForecastAlphaP}{0.107}
\newcommand{\OosRemxChinaForecastCagr}{0.95}
\newcommand{\OosRemxChinaForecastDrawdown}{58.87}
\newcommand{\OosRemxChinaForecastFinal}{10,846}
\newcommand{\OosRemxChinaForecastTime}{52.3}
\newcommand{\OosRemxChinaForecastAlpha}{-2.01}
\newcommand{\OosRemxChinaForecastAlphaP}{0.824}
\newcommand{\OosRemxPassiveCagr}{-4.07}
\newcommand{\OosRemxPassiveDrawdown}{73.34}
\newcommand{\OosRemxPassiveFinal}{6,989}
\newcommand{\OosRemxPassiveTime}{100.0}
\newcommand{\OosRemxPassiveAlpha}{-8.57}
\newcommand{\OosRemxPassiveAlphaP}{0.416}
\newcommand{\OosRemxCashCagr}{2.77}
\newcommand{\OosRemxCashDrawdown}{0.00}
\newcommand{\OosRemxCashFinal}{12,662}
\newcommand{\OosRemxCashTime}{0.0}
\newcommand{\OosLitLowCagr}{4.49}
\newcommand{\OosLitLowDrawdown}{10.68}
\newcommand{\OosLitLowFinal}{14,604}
\newcommand{\OosLitLowTime}{4.0}
\newcommand{\OosLitLowAlpha}{1.61}
\newcommand{\OosLitLowAlphaP}{0.518}
\newcommand{\OosLitHighCagr}{5.40}
\newcommand{\OosLitHighDrawdown}{65.91}
\newcommand{\OosLitHighFinal}{15,736}
\newcommand{\OosLitHighTime}{95.9}
\newcommand{\OosLitHighAlpha}{-1.78}
\newcommand{\OosLitHighAlphaP}{0.819}
\newcommand{\OosLitAvoidCagr}{17.54}
\newcommand{\OosLitAvoidDrawdown}{46.91}
\newcommand{\OosLitAvoidFinal}{40,310}
\newcommand{\OosLitAvoidTime}{55.1}
\newcommand{\OosLitAvoidAlpha}{11.67}
\newcommand{\OosLitAvoidAlphaP}{0.103}
\newcommand{\OosLitMonthlyMACagr}{12.37}
\newcommand{\OosLitMonthlyMADrawdown}{46.06}
\newcommand{\OosLitMonthlyMAFinal}{27,353}
\newcommand{\OosLitMonthlyMATime}{40.6}
\newcommand{\OosLitMonthlyMAAlpha}{5.51}
\newcommand{\OosLitMonthlyMAAlphaP}{0.437}
\newcommand{\OosLitDailyMACagr}{14.04}
\newcommand{\OosLitDailyMADrawdown}{41.13}
\newcommand{\OosLitDailyMAFinal}{31,056}
\newcommand{\OosLitDailyMATime}{41.3}
\newcommand{\OosLitDailyMAAlpha}{9.04}
\newcommand{\OosLitDailyMAAlphaP}{0.204}
\newcommand{\OosLitMomentumCagr}{11.66}
\newcommand{\OosLitMomentumDrawdown}{40.48}
\newcommand{\OosLitMomentumFinal}{25,889}
\newcommand{\OosLitMomentumTime}{47.0}
\newcommand{\OosLitMomentumAlpha}{6.93}
\newcommand{\OosLitMomentumAlphaP}{0.342}
\newcommand{\OosLitMarketCagr}{7.21}
\newcommand{\OosLitMarketDrawdown}{50.49}
\newcommand{\OosLitMarketFinal}{18,237}
\newcommand{\OosLitMarketTime}{56.0}
\newcommand{\OosLitMarketAlpha}{-0.61}
\newcommand{\OosLitMarketAlphaP}{0.931}
\newcommand{\OosLitGlobalForecastCagr}{15.38}
\newcommand{\OosLitGlobalForecastDrawdown}{29.72}
\newcommand{\OosLitGlobalForecastFinal}{34,368}
\newcommand{\OosLitGlobalForecastTime}{50.2}
\newcommand{\OosLitGlobalForecastAlpha}{9.37}
\newcommand{\OosLitGlobalForecastAlphaP}{0.166}
\newcommand{\OosLitChinaForecastCagr}{11.29}
\newcommand{\OosLitChinaForecastDrawdown}{48.85}
\newcommand{\OosLitChinaForecastFinal}{25,163}
\newcommand{\OosLitChinaForecastTime}{51.2}
\newcommand{\OosLitChinaForecastAlpha}{4.11}
\newcommand{\OosLitChinaForecastAlphaP}{0.539}
\newcommand{\OosLitPassiveCagr}{7.33}
\newcommand{\OosLitPassiveDrawdown}{65.91}
\newcommand{\OosLitPassiveFinal}{18,404}
\newcommand{\OosLitPassiveTime}{100.0}
\newcommand{\OosLitPassiveAlpha}{0.07}
\newcommand{\OosLitPassiveAlphaP}{0.993}
\newcommand{\OosLitCashCagr}{2.77}
\newcommand{\OosLitCashDrawdown}{0.00}
\newcommand{\OosLitCashFinal}{12,662}
\newcommand{\OosLitCashTime}{0.0}
\newcommand{\OosCopxLowCagr}{3.26}
\newcommand{\OosCopxLowDrawdown}{16.15}
\newcommand{\OosCopxLowFinal}{13,184}
\newcommand{\OosCopxLowTime}{4.0}
\newcommand{\OosCopxLowAlpha}{0.37}
\newcommand{\OosCopxLowAlphaP}{0.871}
\newcommand{\OosCopxHighCagr}{13.28}
\newcommand{\OosCopxHighDrawdown}{63.48}
\newcommand{\OosCopxHighFinal}{29,314}
\newcommand{\OosCopxHighTime}{95.9}
\newcommand{\OosCopxHighAlpha}{4.30}
\newcommand{\OosCopxHighAlphaP}{0.654}
\newcommand{\OosCopxAvoidCagr}{19.61}
\newcommand{\OosCopxAvoidDrawdown}{43.31}
\newcommand{\OosCopxAvoidFinal}{46,867}
\newcommand{\OosCopxAvoidTime}{55.1}
\newcommand{\OosCopxAvoidAlpha}{13.02}
\newcommand{\OosCopxAvoidAlphaP}{0.126}
\newcommand{\OosCopxMonthlyMACagr}{4.73}
\newcommand{\OosCopxMonthlyMADrawdown}{46.12}
\newcommand{\OosCopxMonthlyMAFinal}{14,897}
\newcommand{\OosCopxMonthlyMATime}{62.5}
\newcommand{\OosCopxMonthlyMAAlpha}{2.25}
\newcommand{\OosCopxMonthlyMAAlphaP}{0.813}
\newcommand{\OosCopxDailyMACagr}{5.40}
\newcommand{\OosCopxDailyMADrawdown}{51.58}
\newcommand{\OosCopxDailyMAFinal}{15,738}
\newcommand{\OosCopxDailyMATime}{60.1}
\newcommand{\OosCopxDailyMAAlpha}{1.83}
\newcommand{\OosCopxDailyMAAlphaP}{0.846}
\newcommand{\OosCopxMomentumCagr}{6.28}
\newcommand{\OosCopxMomentumDrawdown}{44.42}
\newcommand{\OosCopxMomentumFinal}{16,918}
\newcommand{\OosCopxMomentumTime}{62.4}
\newcommand{\OosCopxMomentumAlpha}{2.98}
\newcommand{\OosCopxMomentumAlphaP}{0.753}
\newcommand{\OosCopxMarketCagr}{11.44}
\newcommand{\OosCopxMarketDrawdown}{40.41}
\newcommand{\OosCopxMarketFinal}{25,463}
\newcommand{\OosCopxMarketTime}{47.6}
\newcommand{\OosCopxMarketAlpha}{5.18}
\newcommand{\OosCopxMarketAlphaP}{0.537}
\newcommand{\OosCopxGlobalForecastCagr}{14.35}
\newcommand{\OosCopxGlobalForecastDrawdown}{46.04}
\newcommand{\OosCopxGlobalForecastFinal}{31,799}
\newcommand{\OosCopxGlobalForecastTime}{52.4}
\newcommand{\OosCopxGlobalForecastAlpha}{7.50}
\newcommand{\OosCopxGlobalForecastAlphaP}{0.400}
\newcommand{\OosCopxChinaForecastCagr}{17.19}
\newcommand{\OosCopxChinaForecastDrawdown}{35.60}
\newcommand{\OosCopxChinaForecastFinal}{39,306}
\newcommand{\OosCopxChinaForecastTime}{55.4}
\newcommand{\OosCopxChinaForecastAlpha}{10.12}
\newcommand{\OosCopxChinaForecastAlphaP}{0.249}
\newcommand{\OosCopxPassiveCagr}{13.99}
\newcommand{\OosCopxPassiveDrawdown}{63.62}
\newcommand{\OosCopxPassiveFinal}{30,951}
\newcommand{\OosCopxPassiveTime}{100.0}
\newcommand{\OosCopxPassiveAlpha}{4.91}
\newcommand{\OosCopxPassiveAlphaP}{0.611}
\newcommand{\OosCopxCashCagr}{2.77}
\newcommand{\OosCopxCashDrawdown}{0.00}
\newcommand{\OosCopxCashFinal}{12,662}
\newcommand{\OosCopxCashTime}{0.0}
\newcommand{\OosSpyPassiveCagr}{14.15}
\newcommand{\OosSpyPassiveDrawdown}{33.72}
\newcommand{\OosSpyPassiveFinal}{31,318}
\newcommand{\OosSpyPassiveTime}{100.0}
\newcommand{\OosSpyPassiveAlpha}{-0.09}
\newcommand{\OosSpyPassiveAlphaP}{0.805}
\newcommand{\ArcRemxLowCagr}{7.25}
\newcommand{\ArcRemxLowDrawdown}{26.64}
\newcommand{\ArcRemxLowFinal}{13,844}
\newcommand{\ArcRemxLowTime}{7.0}
\newcommand{\ArcRemxLowAlpha}{3.77}
\newcommand{\ArcRemxLowAlphaP}{0.526}
\newcommand{\ArcRemxHighCagr}{-15.71}
\newcommand{\ArcRemxHighDrawdown}{72.92}
\newcommand{\ArcRemxHighFinal}{4,519}
\newcommand{\ArcRemxHighTime}{92.9}
\newcommand{\ArcRemxHighAlpha}{-21.94}
\newcommand{\ArcRemxHighAlphaP}{0.126}
\newcommand{\ArcRemxAvoidCagr}{-3.63}
\newcommand{\ArcRemxAvoidDrawdown}{60.29}
\newcommand{\ArcRemxAvoidFinal}{8,423}
\newcommand{\ArcRemxAvoidTime}{69.5}
\newcommand{\ArcRemxAvoidAlpha}{-10.03}
\newcommand{\ArcRemxAvoidAlphaP}{0.478}
\newcommand{\ArcRemxMonthlyMACagr}{5.78}
\newcommand{\ArcRemxMonthlyMADrawdown}{41.03}
\newcommand{\ArcRemxMonthlyMAFinal}{12,982}
\newcommand{\ArcRemxMonthlyMATime}{26.8}
\newcommand{\ArcRemxMonthlyMAAlpha}{-1.61}
\newcommand{\ArcRemxMonthlyMAAlphaP}{0.885}
\newcommand{\ArcRemxDailyMACagr}{3.58}
\newcommand{\ArcRemxDailyMADrawdown}{48.91}
\newcommand{\ArcRemxDailyMAFinal}{11,775}
\newcommand{\ArcRemxDailyMATime}{29.8}
\newcommand{\ArcRemxDailyMAAlpha}{-1.94}
\newcommand{\ArcRemxDailyMAAlphaP}{0.871}
\newcommand{\ArcRemxMomentumCagr}{5.31}
\newcommand{\ArcRemxMomentumDrawdown}{44.75}
\newcommand{\ArcRemxMomentumFinal}{12,717}
\newcommand{\ArcRemxMomentumTime}{34.1}
\newcommand{\ArcRemxMomentumAlpha}{-1.79}
\newcommand{\ArcRemxMomentumAlphaP}{0.881}
\newcommand{\ArcRemxMarketCagr}{-8.65}
\newcommand{\ArcRemxMarketDrawdown}{62.38}
\newcommand{\ArcRemxMarketFinal}{6,567}
\newcommand{\ArcRemxMarketTime}{44.9}
\newcommand{\ArcRemxMarketAlpha}{-18.67}
\newcommand{\ArcRemxMarketAlphaP}{0.147}
\newcommand{\ArcRemxGlobalForecastCagr}{4.88}
\newcommand{\ArcRemxGlobalForecastDrawdown}{32.32}
\newcommand{\ArcRemxGlobalForecastFinal}{12,477}
\newcommand{\ArcRemxGlobalForecastTime}{40.7}
\newcommand{\ArcRemxGlobalForecastAlpha}{-2.97}
\newcommand{\ArcRemxGlobalForecastAlphaP}{0.810}
\newcommand{\ArcRemxChinaForecastCagr}{4.58}
\newcommand{\ArcRemxChinaForecastDrawdown}{43.16}
\newcommand{\ArcRemxChinaForecastFinal}{12,311}
\newcommand{\ArcRemxChinaForecastTime}{46.7}
\newcommand{\ArcRemxChinaForecastAlpha}{-4.00}
\newcommand{\ArcRemxChinaForecastAlphaP}{0.759}
\newcommand{\ArcRemxPassiveCagr}{-12.99}
\newcommand{\ArcRemxPassiveDrawdown}{73.34}
\newcommand{\ArcRemxPassiveFinal}{5,238}
\newcommand{\ArcRemxPassiveTime}{99.9}
\newcommand{\ArcRemxPassiveAlpha}{-17.87}
\newcommand{\ArcRemxPassiveAlphaP}{0.242}
\newcommand{\ArcRemxCashCagr}{4.26}
\newcommand{\ArcRemxCashDrawdown}{0.00}
\newcommand{\ArcRemxCashFinal}{12,137}
\newcommand{\ArcRemxCashTime}{0.0}
\newcommand{\ArcLitLowCagr}{6.81}
\newcommand{\ArcLitLowDrawdown}{21.25}
\newcommand{\ArcLitLowFinal}{13,579}
\newcommand{\ArcLitLowTime}{7.0}
\newcommand{\ArcLitLowAlpha}{3.13}
\newcommand{\ArcLitLowAlphaP}{0.545}
\newcommand{\ArcLitHighCagr}{-4.55}
\newcommand{\ArcLitHighDrawdown}{62.30}
\newcommand{\ArcLitHighFinal}{8,054}
\newcommand{\ArcLitHighTime}{92.9}
\newcommand{\ArcLitHighAlpha}{-11.44}
\newcommand{\ArcLitHighAlphaP}{0.276}
\newcommand{\ArcLitAvoidCagr}{3.93}
\newcommand{\ArcLitAvoidDrawdown}{50.96}
\newcommand{\ArcLitAvoidFinal}{11,963}
\newcommand{\ArcLitAvoidTime}{69.5}
\newcommand{\ArcLitAvoidAlpha}{-3.48}
\newcommand{\ArcLitAvoidAlphaP}{0.749}
\newcommand{\ArcLitMonthlyMACagr}{3.92}
\newcommand{\ArcLitMonthlyMADrawdown}{28.94}
\newcommand{\ArcLitMonthlyMAFinal}{11,958}
\newcommand{\ArcLitMonthlyMATime}{34.1}
\newcommand{\ArcLitMonthlyMAAlpha}{-2.96}
\newcommand{\ArcLitMonthlyMAAlphaP}{0.713}
\newcommand{\ArcLitDailyMACagr}{5.05}
\newcommand{\ArcLitDailyMADrawdown}{30.56}
\newcommand{\ArcLitDailyMAFinal}{12,572}
\newcommand{\ArcLitDailyMATime}{34.3}
\newcommand{\ArcLitDailyMAAlpha}{-0.60}
\newcommand{\ArcLitDailyMAAlphaP}{0.943}
\newcommand{\ArcLitMomentumCagr}{13.31}
\newcommand{\ArcLitMomentumDrawdown}{34.50}
\newcommand{\ArcLitMomentumFinal}{17,873}
\newcommand{\ArcLitMomentumTime}{35.4}
\newcommand{\ArcLitMomentumAlpha}{7.97}
\newcommand{\ArcLitMomentumAlphaP}{0.328}
\newcommand{\ArcLitMarketCagr}{7.21}
\newcommand{\ArcLitMarketDrawdown}{34.50}
\newcommand{\ArcLitMarketFinal}{13,820}
\newcommand{\ArcLitMarketTime}{57.2}
\newcommand{\ArcLitMarketAlpha}{0.82}
\newcommand{\ArcLitMarketAlphaP}{0.929}
\newcommand{\ArcLitGlobalForecastCagr}{-1.78}
\newcommand{\ArcLitGlobalForecastDrawdown}{34.50}
\newcommand{\ArcLitGlobalForecastFinal}{9,199}
\newcommand{\ArcLitGlobalForecastTime}{38.5}
\newcommand{\ArcLitGlobalForecastAlpha}{-6.62}
\newcommand{\ArcLitGlobalForecastAlphaP}{0.432}
\newcommand{\ArcLitChinaForecastCagr}{7.18}
\newcommand{\ArcLitChinaForecastDrawdown}{34.50}
\newcommand{\ArcLitChinaForecastFinal}{13,803}
\newcommand{\ArcLitChinaForecastTime}{49.7}
\newcommand{\ArcLitChinaForecastAlpha}{1.86}
\newcommand{\ArcLitChinaForecastAlphaP}{0.839}
\newcommand{\ArcLitPassiveCagr}{-1.88}
\newcommand{\ArcLitPassiveDrawdown}{59.96}
\newcommand{\ArcLitPassiveFinal}{9,155}
\newcommand{\ArcLitPassiveTime}{99.9}
\newcommand{\ArcLitPassiveAlpha}{-8.01}
\newcommand{\ArcLitPassiveAlphaP}{0.484}
\newcommand{\ArcLitCashCagr}{4.26}
\newcommand{\ArcLitCashDrawdown}{0.00}
\newcommand{\ArcLitCashFinal}{12,137}
\newcommand{\ArcLitCashTime}{0.0}
\newcommand{\ArcCopxLowCagr}{9.39}
\newcommand{\ArcCopxLowDrawdown}{7.24}
\newcommand{\ArcCopxLowFinal}{15,171}
\newcommand{\ArcCopxLowTime}{7.0}
\newcommand{\ArcCopxLowAlpha}{5.30}
\newcommand{\ArcCopxLowAlphaP}{0.217}
\newcommand{\ArcCopxHighCagr}{8.73}
\newcommand{\ArcCopxHighDrawdown}{48.55}
\newcommand{\ArcCopxHighFinal}{14,755}
\newcommand{\ArcCopxHighTime}{92.9}
\newcommand{\ArcCopxHighAlpha}{-1.81}
\newcommand{\ArcCopxHighAlphaP}{0.896}
\newcommand{\ArcCopxAvoidCagr}{23.61}
\newcommand{\ArcCopxAvoidDrawdown}{34.94}
\newcommand{\ArcCopxAvoidFinal}{26,769}
\newcommand{\ArcCopxAvoidTime}{69.5}
\newcommand{\ArcCopxAvoidAlpha}{11.71}
\newcommand{\ArcCopxAvoidAlphaP}{0.373}
\newcommand{\ArcCopxMonthlyMACagr}{1.50}
\newcommand{\ArcCopxMonthlyMADrawdown}{50.29}
\newcommand{\ArcCopxMonthlyMAFinal}{10,718}
\newcommand{\ArcCopxMonthlyMATime}{69.7}
\newcommand{\ArcCopxMonthlyMAAlpha}{-5.40}
\newcommand{\ArcCopxMonthlyMAAlphaP}{0.692}
\newcommand{\ArcCopxDailyMACagr}{3.77}
\newcommand{\ArcCopxDailyMADrawdown}{40.75}
\newcommand{\ArcCopxDailyMAFinal}{11,873}
\newcommand{\ArcCopxDailyMATime}{66.2}
\newcommand{\ArcCopxDailyMAAlpha}{-4.39}
\newcommand{\ArcCopxDailyMAAlphaP}{0.744}
\newcommand{\ArcCopxMomentumCagr}{0.12}
\newcommand{\ArcCopxMomentumDrawdown}{47.00}
\newcommand{\ArcCopxMomentumFinal}{10,054}
\newcommand{\ArcCopxMomentumTime}{71.5}
\newcommand{\ArcCopxMomentumAlpha}{-6.37}
\newcommand{\ArcCopxMomentumAlphaP}{0.642}
\newcommand{\ArcCopxMarketCagr}{9.81}
\newcommand{\ArcCopxMarketDrawdown}{35.43}
\newcommand{\ArcCopxMarketFinal}{15,443}
\newcommand{\ArcCopxMarketTime}{53.7}
\newcommand{\ArcCopxMarketAlpha}{3.03}
\newcommand{\ArcCopxMarketAlphaP}{0.796}
\newcommand{\ArcCopxGlobalForecastCagr}{3.68}
\newcommand{\ArcCopxGlobalForecastDrawdown}{43.17}
\newcommand{\ArcCopxGlobalForecastFinal}{11,831}
\newcommand{\ArcCopxGlobalForecastTime}{56.8}
\newcommand{\ArcCopxGlobalForecastAlpha}{-4.48}
\newcommand{\ArcCopxGlobalForecastAlphaP}{0.720}
\newcommand{\ArcCopxChinaForecastCagr}{16.83}
\newcommand{\ArcCopxChinaForecastDrawdown}{35.43}
\newcommand{\ArcCopxChinaForecastFinal}{20,601}
\newcommand{\ArcCopxChinaForecastTime}{64.7}
\newcommand{\ArcCopxChinaForecastAlpha}{6.06}
\newcommand{\ArcCopxChinaForecastAlphaP}{0.627}
\newcommand{\ArcCopxPassiveCagr}{14.48}
\newcommand{\ArcCopxPassiveDrawdown}{42.12}
\newcommand{\ArcCopxPassiveFinal}{18,741}
\newcommand{\ArcCopxPassiveTime}{99.9}
\newcommand{\ArcCopxPassiveAlpha}{3.79}
\newcommand{\ArcCopxPassiveAlphaP}{0.789}
\newcommand{\ArcCopxCashCagr}{4.26}
\newcommand{\ArcCopxCashDrawdown}{0.00}
\newcommand{\ArcCopxCashFinal}{12,137}
\newcommand{\ArcCopxCashTime}{0.0}
\newcommand{\ArcSpyPassiveCagr}{15.21}
\newcommand{\ArcSpyPassiveDrawdown}{22.09}
\newcommand{\ArcSpyPassiveFinal}{19,305}
\newcommand{\ArcSpyPassiveTime}{99.9}
\newcommand{\ArcSpyPassiveAlpha}{-0.15}
\newcommand{\ArcSpyPassiveAlphaP}{0.722}
\newcommand{\fig}[3]{\begin{figure}[!htbp]\centering\includegraphics[width=\textwidth]{figures/#1.pdf}\caption{#2}\label{#3}\end{figure}}
\newcommand{\tabnote}[1]{\par\vspace{4pt}\begin{minipage}{\textwidth}\footnotesize\textit{Notes:} #1\end{minipage}}


\begin{document}
\thispagestyle{plain}
\begin{center}
{\LARGE\bfseries\color{ink} Geopolitical Risk and Returns in Rare\par\vspace{3pt}
Earth and Strategic Metal Markets\par}
\vspace{8pt}
{\large Price responses, investment benchmarks\par\vspace{2pt}
and the information available at the decision date\par}
\end{center}
\vspace{4pt}
\begin{abstract}
\noindent This paper examines whether geopolitical risk explains rare earth prices and adds useful information to investment decisions. Physical-price evidence covers 2010–2026. A chronological forecast exercise begins in 2018, with initial estimation ending in 2017; a separate implementation uses actual archived risk releases from 2022. The selected praseodymium–neodymium (PrNd) model has a positive same-month association with China-related risk changes when previous risk is elevated. Its high-state slope is 3.004 log-return percentage points per standard deviation and remains 3.001 after lagged economic controls. Local projections provide their clearest evidence at the current month, with much wider uncertainty about subsequent price changes. Investment tests compare geopolitical allocations with moving-average rules, time-series momentum, sector rotation and a market-and-trend forecasting model. Several risk filters earn positive historical profits, including \OosLitAvoidCagr\% annually for LIT and \OosCopxAvoidCagr\% for COPX over the revised-data test from 2018. Their advantages do not establish positive factor-adjusted alpha after multiple-testing correction. Adding GPR also fails to improve mean-squared forecasting accuracy consistently. Archived data reveal a separate source of sensitivity: two revised monthly classifications explain a USD 4,706 difference in matched REMX terminal wealth. The evidence supports commodity-specific and state-dependent price associations while making the information requirements for an investable signal explicit.
\end{abstract}
\noindent\small\textbf{Keywords:} geopolitical risk; rare earths; PrNd; strategic metals; return forecasting; investment benchmarks; data vintages.\par
\noindent\small\textbf{JEL classification:} G11; G12; Q02; Q31; F51.\normalsize
\section{Introduction}

Rare earths and strategic metals connect commodity pricing with industrial security. Magnets, batteries, electricity networks and manufacturing equipment depend on access to particular materials and processing stages. Geopolitical tensions can therefore change expected costs and availability. The important empirical question is whether a systematic measure of those tensions explains prices and helps an investor make better decisions.

The supply-chain story does not give a unique investment prediction. A possible shortage can raise the price of a material while increasing the costs faced by firms that use it. A producer may benefit from a higher selling price, but face financing problems or new restrictions of its own. A fund holding companies across several countries and stages of production combines these exposures. The price of a magnet material and the return on a producer-equity portfolio can consequently respond differently to the same news.

The paper asks three questions. First, does a change in geopolitical risk have a different price association when risk was already elevated? Second, does that association remain after accounting for economic activity, the dollar and other commodity prices, and how does it evolve over subsequent months? Third, do geopolitical signals add investment value relative to standard price-based rules and market risk factors? The distinction between explaining a current price change and predicting a future attainable return is central to all three questions.

The main physical series is praseodymium–neodymium mixed metal, abbreviated PrNd. Monthly quote observations cover January 2010–January 2026. The selected interaction model contains 191 returns after allowing for lagged information. Other rare earth quotes, eight IMF material benchmarks and sector equity funds provide additional context. The geopolitical measure is the Caldara–Iacoviello GPR index, particularly its China-related series \citep{caldara2022}.

Investment evaluation uses REMX, LIT and COPX: funds associated with rare earth and strategic materials, the lithium and battery value chain, and copper mining. A chronological exercise starts in March 2018, after the original working paper, with initial estimation ending in December 2017. It uses a consistent revised risk history and is labelled accordingly. A separate 2022–2026 reconstruction reads risk values and refits forecasts from the actual archived release used at each decision. Earlier observations remain useful for estimating historical relationships; their availability as a historical series is distinct from the availability of a live investment signal.

The price results survive a meaningful extension. The selected high-state PrNd slope changes from 3.004 to 3.001 when lagged measures of Chinese activity, global activity, the broad dollar and other commodity returns are included. The evidence for a difference between the two risk states is less precise. Local projections give the clearest positive association in the current month and much wider intervals for later cumulative changes. This pattern is consistent with a contemporaneous price response that need not leave a simple next-month profit opportunity.

The investment results are heterogeneous. Over the revised-data exercise from 2018, avoiding high and rising China-related risk earns \OosLitAvoidCagr\% annually in LIT and \OosCopxAvoidCagr\% in COPX after assumed costs. Standard trend rules also earn positive returns, and conclusions change across funds, execution conventions and calibration windows. Six-factor attribution and nested forecast comparisons do not establish a reliable incremental geopolitical premium. Archived risk data further show how two revised classifications remove losing REMX exposures from a later historical reconstruction.

The central proposition is therefore that \textit{the pricing and investment implications of geopolitical risk are state-dependent, commodity-specific, and sensitive to the information available at the time of the decision}. The contribution is to examine these dimensions together. It combines a transparent conditional price model, economic controls and dynamic responses with benchmarked investment tests and a decision-level account of data revisions. The analysis is exploratory, with its complete specification families and unsuccessful rules retained for replication.

\section{Economic mechanisms and related literature}

A supply channel predicts higher prices when an escalation threatens availability and substitution is difficult. Buyers may accelerate orders or increase precautionary stocks before an actual restriction takes effect. A demand channel can work in the opposite direction: uncertainty can delay industrial investment and reduce expected purchases. Exchange rates, financing conditions and other commodity costs can move with geopolitical news. The sign of an average regression coefficient reflects the combined operation of these forces.

The pre-existing risk environment can matter. During a calm period, a warning may have little effect on purchasing plans. After persistent tension, an additional warning may trigger stockpiling or a revision of quoted availability. An interaction between a new risk change and the previous risk level is a simple way to test this idea. It detects a conditional association; inventories, contracts and lost production are not directly observed here, so it does not identify the supply mechanism by itself.

\citet{caldara2022} establish a news-based framework for measuring adverse geopolitical events and threats. \citet{fan2023} relate geopolitical risk to rare earth trade using Japanese import data, finding higher import unit prices and lower import values, with China-related risk particularly relevant. An import unit value can also change with product composition and trading partners. It therefore motivates the supply-chain hypothesis without fixing the response of a particular quote or equity fund.

\citet{li2023} examine rare earth prices, geopolitical risk and economic activity across time and frequencies. \citet{depraiter2025} study connections with clean-energy investment, and \citet{magnets2025} examine Chinese permanent-magnet exports. The industrial context includes concentrated critical-mineral supply chains and trade restrictions \citep{iea2025}, while regulatory definitions distinguish particular strategic products and uses \citep{eu2024}. These studies support attention to the measured product and its place in the value chain rather than a single coefficient for all materials.

The investment comparison draws on a separate literature. Time-series momentum uses an asset's own past return to determine exposure \citep{moskowitz2012}. Factor models distinguish average returns from compensation associated with familiar equity exposures \citep{fama2015}. Return forecasting must also compete with simple alternatives \citep{campbell2008}. These comparisons make the economic question more demanding: GPR should contribute something beyond a sector's market exposure or a standard price trend, rather than merely improve on one weak passive holding.

\section{Data and the timing of information}
\subsection{Physical materials, traded funds and economic controls}

The main PrNd series is published by SDM Magnetics and attributed by the publisher to Asian Metal assessments \citep{sdm2026}. It is a mischmetal quote in renminbi per metric tonne, not a record of executable transactions. The monthly price is the arithmetic mean of available dated quotes. Missing quote days are not filled before averaging. January 2026 is the final complete month in the retained source; an isolated February observation is excluded.

GUK Magnetics supplies supplementary neodymium and samarium histories from July 2009 to September 2026, described as month-end high prices in USD per kilogram \citep{guk2026}. The IMF workbook provides cobalt, nickel, copper, platinum, palladium, rare earth carbonate, lithium metal and silicon \citep{imfdata,imfmethod}. Five of these series support a January 2011–August 2026 comparison. Rare earth carbonate, lithium metal and silicon begin in June 2012 in the retained workbook and remain explicitly shorter supplementary benchmarks. Carbonate containing 42–45\% rare earth oxides is not an index of all rare earth elements; lithium metal is distinct from lithium carbonate.

\begin{table}[!htbp]
\centering\small
\caption{Data panels and their economic roles}
\label{tab:data}
\begin{tabular}{@{}P{29mm}P{42mm}P{77mm}@{}}
\toprule
Panel & Coverage used & Measurement and role \\
\midrule
PrNd & January 2010–January 2026 & Available publisher quotes; monthly mean; CNY per tonne. Main price model: 191 observations. \\
Other material prices & Long series from 2011; three IMF series from 2012 & IMF benchmarks and separate neodymium/samarium quotes; supplementary material comparison and commodity controls. \\
Sector funds & Daily prices; main test from March 2018 & REMX, LIT and COPX adjusted closes; SPY and EEM as market comparisons or predictors. \\
Economic activity & Monthly; through the PrNd sample where available & China CLI and global real activity; actual China industrial production as a shorter, matched-sample check. \\
Dollar and cash & Dollar from 2006; cash over all test dates & Broad dollar changes for the price model; daily effective federal funds rate for cash accrual. \\
Equity factors & Daily; through September 2026 & US five-factor returns and momentum; ex-post risk adjustment. \\
\bottomrule
\end{tabular}
\end{table}

Daily adjusted prices for REMX, LIT, COPX and SPY are retained from the market-data provider; EEM adds an emerging-market predictor \citep{yahoo2026}. Fund definitions and inception dates are checked against issuer documentation \citep{vaneck2026,globalxlit,globalxcopx}. These are equity investments. REMX holds firms involved in the strategic-material value chain; LIT includes mining, refining and battery-related production; COPX represents copper-mining firms. Their historical distributions and corporate actions follow the provider's adjustment methodology. Fund expenses already affect prices and are not deducted again.

The economic controls add four dimensions. China's amplitude-adjusted composite leading indicator, or CLI, measures its business-cycle environment. The Dallas Fed global real activity index is based on bulk dry cargo shipping rates and serves as a global commodity-demand proxy \citep{dallas2026}. Broad dollar log changes capture a financial and exchange-rate influence. An equal-weight average of copper and Brent log returns captures a wider commodity cycle. PrNd remains in CNY; the dollar control does not mechanically convert its quote into dollars.

Actual Chinese industrial production is examined separately using the OECD index relative to the same month of the previous year \citep{oecd2026}. This source ends in November 2023 and has missing January observations in several years. They are not imputed. With a one-month control lag, the resulting price model ends in December 2023 and contains 152 complete observations. An uncontrolled model is re-estimated on exactly those dates, so a change in the coefficient is not automatically attributed to the added control when the sample also changes. The longer CLI is labelled as a leading indicator, not as industrial production.

\subsection{Historical observations and historical availability}

The original GPR working paper was publicly available in 2018 \citep{caldara2018}. This makes 2018 a reasonable conservative boundary for a post-publication investment exercise. It does not establish that every current country series, methodology or historical value was available then. The older GPR page describes monthly updates around the tenth calendar day; the methodology changed in October 2021 \citep{gprold2021}. Current documentation describes early-month updates and possible revisions \citep{gprdata}.

The distinction affects the design. The revised-data exercise starts on 15 March 2018 and trades at the eleventh observed US trading-session close of each month, allowing time after the older update convention. All monthly price-based comparisons use the same execution dates. It uses the November 2026 risk history and is a retrospective chronological test, rather than a reconstruction of every original 2018 release. Initial model estimation ends in 2017, but the evaluation period was examined during research; it is not an untouched confirmation sample.

The archive exercise instead starts on 8 March 2022, uses 56 retained monthly releases, and trades at the sixth session under the later early-month convention. Forecast coefficients are refitted from the historical risk data in each release, not from the 2026 replacement for those data. Its threshold is fixed from the 2009–2018 history in a January 2022 archive file. Exact historical server upload timestamps are not audited; availability is inferred from the named release and documented update convention.

\begin{table}[!htbp]
\centering\small
\caption{Three distinct uses of historical information}
\label{tab:timing}
\begin{tabular}{@{}P{35mm}P{50mm}P{63mm}@{}}
\toprule
Exercise & Risk calibration or version & Interpretation \\
\midrule
Price associations, 2010–2026 & Selected 2009–2022 state definition; latest data & Retrospective explanation of historical quote returns. \\
Chronological test from 2018 & Thresholds from 2009–2017; initial forecasts estimated through 2017; revised risk history & Sequential forecasts using earlier closed outcomes, with explicit limits on original data availability and model selection. \\
Archived test, 2022–2026 & January 2022 frozen threshold; each month's own release for signals and refitting & Reconstruction using retained risk versions, conditional on the stated release convention. \\
\bottomrule
\end{tabular}
\end{table}

The retained old-method workbook supplies an additional 2018–2021 sensitivity check. It contains a frozen 2021 history, including a differently scaled China series. Thresholds are recalibrated within that method before comparing classifications. This snapshot is useful evidence of method sensitivity, but is not presented as a complete sequence of 2018 vintages.

Cash accrues the previous calendar day's effective federal funds rate using daily ACT/360 compounding, including weekends between trading closes \citep{fred2026}. It is an opportunity-cost proxy; an individual broker rate can differ. All results are before tax and in USD. Factor attribution uses the risk-free series supplied with the French factors to match their definitions \citep{french2026}; this is distinct from the cash assumption used in portfolio allocation.

\fig{01_prices_and_risk}{PrNd prices and China-related geopolitical risk. The price is a monthly mean of available publisher quotes, while risk is a news-based measure in original units. The panels have different units. A high price level is different from a high monthly return.}{fig:prices}

\section{Price models and dynamic responses}

The monthly quote return is
\begin{equation}
r_t=100\,[\ln(P_t)-\ln(P_{t-1})].
\label{eq:return}
\end{equation}
Here $P_t$ is the monthly price and $\ln$ denotes the natural logarithm. A price increase from 100 to 110 is a 10\% simple return and about 9.53 log-return percentage points. Regressions use log returns; portfolio wealth uses compounded simple returns.

Risk enters as $g_t=\ln(1+\mathrm{GPR}_t)$. Its level describes how elevated the measured risk is; $\Delta g_t=g_t-g_{t-1}$ describes whether it increased. Adding one allows very small or zero country values. For the selected price model, $z_t$ is the risk change centred and divided by its standard deviation over 2009–2022. A one-unit increase in $z_t$ is therefore a one-standard-deviation risk change, rather than an arbitrary unit of the source index.

The state-dependent regression is
\begin{equation}
r_t=\alpha+\beta z_t+\gamma H_{t-1}+\delta z_tH_{t-1}
       +\phi r_{t-1}+\symbf{\theta}'\symbf{x}_{t-1}+u_t.
\label{eq:state}
\end{equation}
The constant $\alpha$ allows for average return, and $r_{t-1}$ allows for persistence in the price's own movement. The indicator $H_{t-1}$ equals one when previous-month China-related risk exceeds the 80th percentile of its 2009–2022 calibration history, and zero otherwise. An 80th percentile is exceeded by roughly one fifth of the calibration observations; it need not classify one fifth of a later sample as high risk. The vector $\symbf{x}_{t-1}$ contains the lagged economic controls. The unexplained component is $u_t$.

The product $z_tH_{t-1}$ is an interaction. In the lower state, $H_{t-1}=0$ and the slope on a risk change is $\beta$. In the high state it is $\beta+\delta$. The separate state term $\gamma H_{t-1}$ also allows average return to differ. For example, a high-state slope of three means about three log-return percentage points associated with a typical risk increase, holding the other terms fixed. It does not mean that every high-risk month earns three percent.

The original model sets the economic-control coefficients to zero. The primary extension adds lagged China CLI, global activity, dollar returns and the copper–Brent basket return. A further check uses their same-month values; this answers a more conditional question because contemporaneous controls may themselves respond to the same information. The industrial-production model uses its actual shorter coverage and a matched baseline. These are robustness exercises for the selected association, not newly independent discoveries.

Dynamic responses are estimated with local projections \citep{jorda2005}. At each horizon $h$, the dependent variable is the cumulative price change
\begin{equation}
y_{t,h}=100\,[\ln(P_{t+h})-\ln(P_{t-1})],\qquad h=0,\ldots,12.
\label{eq:lp}
\end{equation}
The regression retains the current risk change, prior state and their interaction, adds three lags each of returns and risk changes, and uses the lagged economic controls. At $h=0$, the outcome is the current monthly return; at $h=3$, it is the cumulative change through three further months. Estimating a separate regression at each horizon lets the data describe the timing without imposing a recursive VAR response.

The main response plot uses the same 177 dates at every horizon, from May 2010 to January 2025; all required later prices are observed by January 2026. Horizon-specific samples are retained separately. The high-state response is again the sum of the risk and interaction slopes. It is a response conditional on the included variables, not the effect of an externally identified geopolitical shock.

Uncertainty uses Bartlett HAC covariance estimates with a finite-sample correction, following \citet{newey1987}. The main regressions allow three monthly lags; local projections use the larger of three and $h+1$. Scores are placed back on the monthly calendar before covariance calculation, preserving gaps in the industrial-production sample. Reported intervals are pointwise 95\% bands. The wider original search and its Holm and false-discovery-rate adjustments remain in the archive \citep{holm1979,benjamini1995}. Nominal significance in a selected model is interpreted alongside that selection history.

\section{Price evidence after economic controls}

The selected original PrNd model contains 191 observations, including 63 months following elevated risk. Its high-state slope is 3.004 log-return percentage points per standard deviation, with nominal $p=0.019$. The lower-state slope is 0.160 and is estimated imprecisely. The estimated difference is 2.844, with nominal $p=0.044$. Three log-return percentage points correspond to a simple price change of approximately $100[\exp(3.004/100)-1]=3.05\%$ for that marginal shift, holding the other terms fixed.

Table~\ref{tab:controls} shows that adding lagged economic controls leaves the high-state estimate almost unchanged at 3.001, with nominal $p=0.023$. The same-month-control version gives 2.790, with $p=0.041$. These comparisons weaken a simple explanation based only on the included business-cycle, dollar and commodity variables. They do not eliminate all omitted influences or identify an exogenous geopolitical effect. The difference between the high- and lower-state slopes is also less precise after controls: its $p$-value is 0.053 with lagged controls and 0.080 with current controls.

\begin{table}[!htbp]
\centering\small
\caption{PrNd high-state associations with economic controls}
\label{tab:controls}
\setlength{\tabcolsep}{4pt}
\begin{tabular}{@{}lrrrrr@{}}
\toprule
Model & $N$ & High slope & SE & $p$ high & $p$ difference \\
\midrule
Original model & 191 & 3.004 & 1.284 & 0.019 & 0.044 \\
Lagged economic controls & 191 & 3.001 & 1.321 & 0.023 & 0.053 \\
Same-month economic controls & 191 & 2.790 & 1.367 & 0.041 & 0.080 \\
Original: IP-matched dates & 152 & 5.199 & 1.905 & 0.006 & 0.009 \\
Lagged IP controls & 152 & 5.347 & 1.935 & 0.006 & 0.010 \\
\bottomrule
\end{tabular}
\tabnote{The risk scaling and prior-state threshold are held fixed across these comparisons. The first three rows use 191 observations. The industrial-production (IP) rows use the same 152 complete dates. The last column tests the difference from the lower-state slope. Nominal HAC inference is conditional on the specification; coefficients are log-return percentage points per standard deviation of the risk change.}
\end{table}

The industrial-production check illustrates why matched samples are necessary. With lagged IP and the other economic controls, the high-state estimate is 5.347. The original model on those same dates gives 5.198. Much of the increase from the full-sample 3.004 therefore reflects which observations remain, rather than an effect created by adding industrial production. The IP series has not been extended beyond its genuine coverage or interpolated across missing months.

\fig{02_price_controls}{High-state PrNd coefficients and pointwise 95\% HAC intervals. The original and lagged-control full-sample estimates are nearly identical. The two industrial-production comparisons use a common shorter sample and should be compared with each other. Positive conditional price associations do not establish an identified causal mechanism.}{fig:controls}

Other earlier checks remain informative. Deleting individual calendar years leaves the selected high-state slope positive, between about 2.12 and 3.64. Six- and twelve-month-block bootstrap intervals are positive for the high-state slope, while evidence for the difference between states is less uniform. An alternative calibration fixed entirely from 2000–2009 gives a high-state slope of 1.111, with nominal $p=0.074$, and $p=0.461$ for the state difference. That alternative changes both the standardisation and the number of months classified as high risk. The current controls do not resolve this dependence on the definition of the state.

The broader specification search also limits discovery claims. The original archive retains 256 linear specifications, 30 quantile specifications, ten state specifications and twelve daily comparisons, together with later and earlier subsamples. A simple ten-high-slope Bonferroni correction changes the selected model's 0.019 nominal $p$-value to about 0.193. The new economic controls test the robustness of that selected association; they do not provide a separate, previously untouched confirmation.

\subsection{When the conditional response occurs}

The controlled local projections provide their clearest evidence at the current month. On the common 177-date sample, the high-state $h=0$ estimate is 3.118, with a pointwise interval of approximately $[0.201,6.036]$. It differs slightly from Table~\ref{tab:controls} because the dynamic regression includes additional return and risk lags and uses common dates suitable for every horizon.

At one further month, the cumulative estimate is 2.744, with an interval of about $[-3.304,8.793]$. At six months it is 1.354, with $[-5.717,8.425]$, and at twelve months it is 0.299, with $[-7.582,8.180]$. The point estimates fade, but the wider intervals do not establish a precisely measured reversal or zero long-run effect. They show that these observations provide substantially less information about a persistent response than about the current-month association.

\fig{03_local_projections}{Controlled cumulative PrNd associations after a risk change in the high prior-risk state. Each point is estimated in a separate local projection using the same 177 dates, three return lags, three risk-change lags and lagged economic controls. The shaded region gives pointwise 95\% HAC intervals; it is not a simultaneous confidence band. The measured change is not an externally identified shock.}{fig:lp}

The economic implication is important for the investment test. A quote can move while geopolitical news is arriving during the month. Someone acting on the completed monthly index then faces a later return, in a different traded vehicle. A substantial current-month coefficient is therefore compatible with limited subsequent forecasting value. The wider material comparison also remains heterogeneous: the original extended mean-return screen finds no common significant lagged-risk relation across the material universe. Absolute-return and quantile results describe different parts of the distribution and are not treated as automatic buy signals \citep{koenker1978}.

\section{Investment design and competing signals}
\subsection{Allocation rules and transaction timing}

Each main strategy starts with USD 10,000 and alternates between a full position in one ETF and cash. There is no leverage or short selling. At a monthly decision, it uses the completed previous month's risk or price information. The standard revised-data test executes at the eleventh trading-session close from 15 March 2018; the archive reconstruction uses the sixth session from 8 March 2022. A newly selected position applies after that close, so it does not earn the movement that has already occurred before execution.

For the 2018 exercise, high risk is defined by the 80th percentile of the 2009–2017 history. Six simple geopolitical rules are retained in each fund: hold when China-related risk is low; hold when it is high; hold when it is both high and rising; hold except when it is high and rising; hold when global risk is high and rising; or hold whenever China-related risk is rising. Level and change are separate conditions. Risk can be elevated and falling, or low and rising. These six interpretations are evaluated together rather than reporting only a favourable variant.

The benchmarks ask whether GPR adds something beyond familiar price information. The monthly MA200 rule compares the last completed month's closing price with its 200-session moving average and uses the same monthly execution date as GPR. The daily MA200 rule reads that comparison at the preceding trading close and executes at the next close. This explicit one-session delay avoids using the new closing price to claim a position that already earned that day's return.

The time-series momentum benchmark holds the fund when its completed trailing twelve-month return is positive, and otherwise holds cash. A sector-rotation benchmark ranks REMX, LIT and COPX on their twelve-minus-one-month return: it excludes the most recent month, selects the highest positive past return and otherwise holds cash. This is a three-fund rotation comparison, not a claim to represent momentum across the whole stock market. SPY provides a broad equity opportunity-cost comparison.

Costs are ten basis points, or 0.1\%, per one-way unit of portfolio turnover. An initial USD 10,000 purchase therefore costs about USD 10, and a later exit incurs another charge at the then-current value. Unchanged weights incur no additional turnover cost. Moving from one fully invested ETF to another involves a sale and a purchase. Initial purchases and final liquidations are charged, and higher-cost alternatives are retained.

The wealth calculation is
\begin{equation}
W_d=W_{d-1}\bigl[1+w_{d-1}R^E_d+(1-w_{d-1})R^c_d\bigr]
                  \bigl[1-c|w_d-w_{d-1}|\bigr].
\label{eq:wealth}
\end{equation}
Here $W_d$ is wealth after the decision at trading close $d$, $w_d$ is the ETF weight of zero or one, $R^E_d$ is the adjusted return between consecutive closes, $R^c_d$ is cash income over that interval, and $c$ is the one-way proportional cost. The old weight earns the interval's return, followed by the cost of changing the position. The rotation benchmark uses the corresponding vector of weights and the sum of absolute weight changes.

\subsection{Forecasting beyond a market-and-trend model}

A separate nested comparison evaluates incremental predictive information. The benchmark forecasts the ETF return to the next monthly execution close using five already observed price variables: the previous calendar month's own return, twelve-minus-one-month momentum, the distance from MA200, the SPY return and the EEM return. An augmented version adds the prior completed month's global GPR level, change, high-state indicator and interaction. Another adds the analogous China-related terms. The two risk additions are evaluated separately.

Initial estimation uses only return labels ending by December 2017, with 73 usable examples at the first 2018 forecast. Thereafter the window expands. Every training label must end strictly before the current execution close, so a return ending at today's close cannot enter a model used to trade at that same close. A minimum of sixty labels is required. The strategy holds the ETF only when its predicted return exceeds a cash forecast based on the previously observed rate.

For the archive exercise, historical risk features and coefficients are rebuilt from each month's actual release. Price features remain based on already completed months, and the same strict endpoint rule applies. This is stronger than simply replacing the last risk observation in a model trained on the 2026 history. All forecast origins, training counts and last training-label dates are supplied in the replication files.

Forecast accuracy is measured relative to the market-and-trend benchmark:
\begin{equation}
R^2_{\mathrm{forecast}}
=1-\frac{\sum_t(R_t-\widehat R^{\mathrm{GPR}}_t)^2}
         {\sum_t(R_t-\widehat R^{\mathrm{base}}_t)^2}.
\label{eq:r2}
\end{equation}
A positive 2\% means that adding GPR reduces mean-squared prediction error by 2\%; a negative value means a larger error. The Clark–West comparison adjusts the nested-model error difference for estimation noise from additional parameters \citep{clark2007}. One-sided HAC $p$-values and a Holm correction across the six fund-by-risk additions are reported. These remain retrospective comparisons, given the broader research and data-vintage limitations.

\section{Results against investment benchmarks}
\subsection{Historical profit and the price-based alternatives}

Table~\ref{tab:benchmark} gives the common-date 2018–2026 results after costs. The threshold calibration, execution delay and period are stated together because all affect the outcome. CAGR is the constant annual growth rate that reproduces the final compounded wealth; it is not the arithmetic mean of yearly returns. Maximum drawdown is the largest decline from a preceding portfolio peak.

\begin{table}[!htbp]
\centering\small
\caption{Net annualised returns, 15 March 2018–30 October 2026}
\label{tab:benchmark}
\begin{tabular}{@{}lrrr@{}}
\toprule
Rule & REMX & LIT & COPX \\
\midrule
Low China risk & 6.87 & 4.49 & 3.26 \\
High China risk & -7.89 & 5.40 & 13.28 \\
High and rising China risk & -12.90 & -7.39 & -3.34 \\
Avoid high and rising China risk & 11.71 & 17.54 & 19.61 \\
High and rising global risk & -11.31 & -6.91 & -3.66 \\
China risk rising & -12.90 & -7.39 & -3.34 \\
MA200: monthly decisions & 10.99 & 12.37 & 4.73 \\
MA200: daily decisions & 11.27 & 14.04 & 5.40 \\
12-month momentum & 7.85 & 11.66 & 6.28 \\
Market-and-trend forecast & 1.68 & 7.21 & 11.44 \\
Passive fund & -4.07 & 7.33 & 13.99 \\
Cash proxy & 2.77 & 2.77 & 2.77 \\
\bottomrule
\end{tabular}
\tabnote{Entries are CAGRs in percent. Monthly decisions execute at the eleventh session; daily MA200 uses a prior-close signal and next-close execution. The GPR history is revised data, with thresholds calibrated from 2009–2017. Costs are ten basis points per one-way turnover. The full output additionally reports drawdowns, Sharpe ratios, exposure, costs and every forecast-based allocation.}
\end{table}

The rule avoiding high and rising China-related risk earns \OosLitAvoidCagr\% annually in LIT and \OosCopxAvoidCagr\% in COPX. Passive returns are \OosLitPassiveCagr\% and \OosCopxPassiveCagr\%, respectively. These rules also exceed the daily MA200 returns of \OosLitDailyMACagr\% and \OosCopxDailyMACagr\%, and the twelve-month momentum returns of \OosLitMomentumCagr\% and \OosCopxMomentumCagr\%. The comparison provides a concrete reason to investigate the geopolitical filter beyond its improvement on passive sector ownership.

The risk experience still matters. The LIT candidate has a \OosLitAvoidDrawdown\% maximum drawdown, compared with \OosLitDailyMADrawdown\% for daily MA200 and \OosLitPassiveDrawdown\% passively. COPX's candidate drawdown is \OosCopxAvoidDrawdown\%, versus \OosCopxDailyMADrawdown\% for daily MA200 and \OosCopxPassiveDrawdown\% passively. A higher final CAGR can therefore come with a different and sometimes less favourable path of losses.

\fig{04_benchmark_wealth}{Geopolitical filters and standard alternatives on common 2018–2026 dates. The GPR candidate holds LIT or COPX except after high and rising China-related risk. Daily MA200, passive fund ownership and the cash proxy share the same initial capital and costs. The logarithmic scale compares proportional wealth changes. These are retrospective results using revised GPR data.}{fig:wealth}

REMX gives a different comparison. Its low-China-risk rule earns \OosRemxLowCagr\%, against \OosRemxDailyMACagr\% for daily MA200 and \OosRemxMomentumCagr\% for momentum. Its \OosRemxLowDrawdown\% drawdown is substantially smaller, so the return ranking is not a complete investor-utility comparison. The simple high-and-rising global-risk rule loses money. Positive PrNd price evidence has not translated into a universal directional rare earth equity rule.

The three-fund rotation benchmark earns about $-5.12\%$ annually, while passive SPY earns \OosSpyPassiveCagr\% with a \OosSpyPassiveDrawdown\% maximum drawdown. The rotation result is retained despite being unfavourable. SPY is a different exposure, but it makes the opportunity cost of the sector allocation visible rather than treating improvement on a sector's own passive history as sufficient investment evidence.

\subsection{Factor-adjusted performance and uncertainty}

To separate average returns from familiar equity exposures, daily excess portfolio returns are regressed on the US market, size, value, profitability, investment and momentum factors \citep{fama2015,french2026}:
\begin{equation}
R_{p,d}-R_{f,d}=\alpha+\symbf{b}'\symbf{F}_d+e_d.
\label{eq:alpha}
\end{equation}
Here $R_{f,d}$ is the risk-free return supplied with the factors, $\symbf{F}_d$ is the six-factor vector, and the intercept $\alpha$ is the average return left after those static exposures. Annualised alpha is $252\alpha$ and is an arithmetic rate, not CAGR. The opening transaction cost is included in the first holding-interval return. Five-lag HAC inference is used. The downloaded factors end in September 2026, so attribution ends then, while portfolio wealth continues through October. US factors do not fully span every foreign-country and commodity exposure of these international funds.

\begin{table}[!htbp]
\centering\small
\caption{Selected geopolitical allocations and factor-adjusted comparisons}
\label{tab:alpha}
\setlength{\tabcolsep}{4pt}
\begin{tabular}{@{}llrrrr@{}}
\toprule
ETF & Rule & CAGR & Max. DD & Alpha & $p_{\alpha}$ \\
\midrule
REMX & GPR candidate & 6.87 & -12.16 & 4.01 & 0.176 \\
REMX & MA200: daily decisions & 11.27 & -52.01 & 8.41 & 0.342 \\
REMX & 12-month momentum & 7.85 & -49.16 & 6.74 & 0.456 \\
LIT & GPR candidate & 17.54 & -46.91 & 11.67 & 0.103 \\
LIT & MA200: daily decisions & 14.04 & -41.13 & 9.04 & 0.204 \\
LIT & 12-month momentum & 11.66 & -40.48 & 6.93 & 0.342 \\
COPX & GPR candidate & 19.61 & -43.31 & 13.02 & 0.126 \\
COPX & MA200: daily decisions & 5.40 & -51.58 & 1.83 & 0.846 \\
COPX & 12-month momentum & 6.28 & -44.42 & 2.98 & 0.753 \\
\bottomrule
\end{tabular}
\tabnote{The GPR rows are REMX low-China-risk and the LIT/COPX rules avoiding high and rising China risk. CAGR and drawdown cover the full 2018–2026 portfolio window. Alpha is the annualised six-factor intercept through September 2026. All rates are percentages. Displayed $p$-values are nominal; none of these GPR alpha results is significant after adjustment across the 24 GPR-related candidates.}
\end{table}

The LIT avoidance candidate has an annualised alpha of approximately \OosLitAvoidAlpha\%, with nominal $p=\OosLitAvoidAlphaP$; COPX's is about \OosCopxAvoidAlpha\%, with $p=\OosCopxAvoidAlphaP$. Their confidence intervals remain wide. The low-risk REMX alpha is approximately \OosRemxLowAlpha\%, with $p=\OosRemxLowAlphaP$. Positive point estimates therefore remain candidates for further evaluation, rather than established residual premiums.

\fig{05_factor_alpha}{Annualised six-factor alphas and pointwise 95\% HAC intervals. A positive point estimate with an interval crossing zero is imprecise evidence of performance beyond the included factors. Daily MA200 and twelve-month momentum are shown alongside the GPR candidates. The factor sample ends in September 2026 and the factors are US equity benchmarks.}{fig:alpha}

For each fund, paired monthly log returns also compare GPR allocations directly with cash, passive ownership, monthly MA200, daily MA200, momentum and the market-and-trend forecast allocation. Each comparison family applies Holm adjustment across eighteen direct GPR rules and six GPR-augmented forecast allocations. None establishes positive outperformance at a 5\% adjusted level. Block-bootstrap intervals preserve six- or twelve-month groups \citep{kunsch1989}. For example, the six-month-block interval for LIT's annualised mean log advantage over monthly MA200 is approximately $[-12.35,24.13]$ percentage points; COPX's is $[-5.75,32.05]$. These are not confidence intervals for CAGR.

\fig{06_annual_returns_clear}{Calendar-year performance with separate, horizontal year labels. Cell values are percent returns after costs, with the colour scale capped at $\pm50\%$. The first year begins at the March 2018 execution date and 2026 ends in October, so those two cells are partial calendar years. The complete year-by-year benchmark table is supplied separately.}{fig:annual}

\subsection{Does adding GPR improve a return forecast?}

The portfolio results and forecast-error results answer different questions. In the revised-data test, all six global- or China-risk additions have negative forecast $R^2$ relative to the market-and-trend model. The range is approximately $-1.44\%$ to $-8.59\%$. Clark–West nominal one-sided $p$-values are about 0.17–0.31 and become one after the six-comparison correction. The data do not establish an improvement in mean-squared accuracy from these added risk variables.

\begin{table}[!htbp]
\centering\footnotesize
\setlength{\tabcolsep}{3.5pt}
\caption{Incremental forecast accuracy and the corresponding allocations}
\label{tab:forecast}
\begin{tabular}{@{}lllrrrrr@{}}
\toprule
Sample & ETF & Added GPR & $N$ & $R^2$ & CAGR & $p_{CW}$ & $p_{Holm}$ \\
\midrule
2018–2026 & REMX & Global & 103 & -5.59 & 16.15 & 0.190 & 1.000 \\
2018–2026 & REMX & China & 103 & -1.46 & 0.95 & 0.308 & 1.000 \\
2018–2026 & LIT & Global & 103 & -5.26 & 15.38 & 0.180 & 1.000 \\
2018–2026 & LIT & China & 103 & -2.22 & 11.29 & 0.250 & 1.000 \\
2018–2026 & COPX & Global & 103 & -8.59 & 14.35 & 0.171 & 1.000 \\
2018–2026 & COPX & China & 103 & -1.44 & 17.19 & 0.190 & 1.000 \\
2022–2026 & REMX & Global & 55 & -6.93 & 4.88 & 0.476 & 1.000 \\
2022–2026 & REMX & China & 55 & -0.94 & 4.58 & 0.442 & 1.000 \\
2022–2026 & LIT & Global & 55 & -8.80 & -1.78 & 0.518 & 1.000 \\
2022–2026 & LIT & China & 55 & -0.48 & 7.18 & 0.316 & 1.000 \\
2022–2026 & COPX & Global & 55 & -17.39 & 3.68 & 0.565 & 1.000 \\
2022–2026 & COPX & China & 55 & 2.71 & 16.83 & 0.115 & 0.687 \\
\bottomrule
\end{tabular}
\tabnote{$R^2$ is a percentage relative to the same market-and-trend forecast; negative values mean larger error. $N$ counts completed holding periods: 103 from 2018 and 55 in the archive exercise. The final incomplete forecast period is excluded from accuracy but its decision remains in the wealth simulation. CAGR is the return from converting the respective forecast to an ETF/cash decision. CW denotes the one-sided Clark–West comparison; Holm adjustment covers six additions per sample.}
\end{table}

A forecast allocation can still earn money without reducing squared error. It trades on whether a prediction exceeds the cash threshold, not on the full distance between prediction and realised return. For example, the REMX global-GPR forecast allocation earns \OosRemxGlobalForecastCagr\% annually, compared with \OosRemxMarketCagr\% for the price-only forecast allocation, while its squared prediction error is 5.59\% larger. This is not a contradiction: the investment rule and the forecast loss function reward different features of the predictions. It also makes the particular profitable allocation an insufficient basis for claiming general forecasting superiority.

\fig{07_incremental_forecasts}{Forecast accuracy added by global or China-related GPR. The benchmark already includes own-return dynamics, momentum, a moving-average gap and broad-market returns. Negative bars mean larger mean-squared errors. The left panel uses revised historical risk data from 2018; the right refits models from each actual archived risk release from 2022.}{fig:forecasts}

\section{Archived releases, calibration and revision attribution}
\subsection{Refitting from the information in each release}

The archive exercise repeats the comparisons on 8 March 2022–30 October 2026. Each month's geopolitical allocation and risk-augmented forecast uses that release's risk history, with the threshold fixed from the January 2022 archive. The nested models are re-estimated using only return labels that ended before the current decision close. Standard price-based rules are evaluated on the same dates. This directly addresses the difference between using historically available risk values and a later revised historical file.

\begin{table}[!htbp]
\centering\small
\caption{Net annualised returns using archived GPR releases, 2022–2026}
\label{tab:archive}
\begin{tabular}{@{}lrrr@{}}
\toprule
Rule & REMX & LIT & COPX \\
\midrule
Low China risk & 7.25 & 6.81 & 9.39 \\
High China risk & -15.71 & -4.55 & 8.73 \\
High and rising China risk & -6.84 & -2.59 & -4.44 \\
Avoid high and rising China risk & -3.63 & 3.93 & 23.61 \\
High and rising global risk & 7.57 & 7.53 & 16.02 \\
China risk rising & -6.84 & -2.59 & -4.44 \\
MA200: monthly decisions & 5.78 & 3.92 & 1.50 \\
MA200: daily decisions & 3.58 & 5.05 & 3.77 \\
12-month momentum & 5.31 & 13.31 & 0.12 \\
Market-and-trend forecast & -8.65 & 7.21 & 9.81 \\
Market model + global GPR & 4.88 & -1.78 & 3.68 \\
Market model + China GPR & 4.58 & 7.18 & 16.83 \\
Passive fund & -12.99 & -1.88 & 14.48 \\
Cash proxy & 4.26 & 4.26 & 4.26 \\
\bottomrule
\end{tabular}
\tabnote{Entries are CAGRs in percent. Monthly execution is at the sixth trading-session close; daily MA200 retains its one-session delay. Thresholds use the 2009–2018 history in the January 2022 archive. Risk-augmented models are refitted from each release's historical values. Costs are ten basis points one way. This sample differs from the 2018 exercise in dates, calibration and execution convention, so differences between the tables are not attributed solely to revisions.}
\end{table}

The results again depend on the vehicle. REMX's low-risk rule earns \ArcRemxLowCagr\% annually, compared with \ArcRemxDailyMACagr\% for daily MA200, \ArcRemxMomentumCagr\% for momentum and \ArcRemxCashCagr\% for cash. It is invested during only four monthly intervals. LIT's avoidance rule earns \ArcLitAvoidCagr\%, while its momentum benchmark earns \ArcLitMomentumCagr\%. COPX's avoidance rule earns \ArcCopxAvoidCagr\%, against \ArcCopxPassiveCagr\% passively and \ArcCopxDailyMACagr\% for daily MA200. A result that is attractive in one fund and period therefore does not carry over automatically to another.

The forecast comparison gives a small positive accuracy result for the archived COPX China-risk addition: its $R^2$ relative to the market-and-trend model is 2.71\%. The nominal Clark–West $p$-value is 0.115 and the six-comparison-adjusted value is 0.687. Other archive risk additions have negative error-based $R^2$. This is a useful candidate for further observation, with limited precision over 55 completed holding periods, rather than a broad validation of geopolitical forecasting.

\subsection{The threshold remains part of the economic result}

Fixed and adaptive thresholds are compared without selecting the best percentile after observing performance. The alternatives are the 70th, 80th and 90th percentiles of either the fixed earlier history or the preceding 120 months. The adaptive reference excludes the month being classified. In the archive version, its historical values come from the current archived release, not the final revised history.

For the 2018–2026 avoidance rule at the 80th percentile, LIT earns 17.54\% under the fixed reference and 20.32\% under the adaptive reference; COPX earns 19.61\% and 25.76\%. In the archived 2022–2026 version, LIT changes from 3.93\% fixed to approximately $-0.04\%$ adaptive, while COPX changes from 23.61\% to 22.66\%. These comparisons show why a favourable implementation should remain attached to its precise state definition and sample.

\fig{08_threshold_stability}{Fixed and adaptive thresholds for the rule avoiding high and rising China-related risk. The upper panels use the revised-data test from 2018; the lower panels use archived releases from 2022. Adaptive thresholds use the previous 120 months and exclude the classified month. Each panel retains all three tested percentiles. The panels differ in evaluation period and execution convention.}{fig:thresholds}

The long-history diagnostic also remains relevant. Over 2011–2026, a COPX high-risk rule with a fixed 2000–2009 threshold earns 7.70\%, but its trailing-120-month version earns $-1.37\%$. Those complete results are retained in the supplementary files as retrospective context, not evidence of a live 2011 signal. The sensitivity raises a substantive question: a fixed risk threshold may partly identify persistent shifts in the risk distribution and market cycle. The added benchmarks and macro controls examine that concern but do not make it disappear.

\subsection{Which revised decisions change portfolio wealth?}

A matched comparison isolates the effect of numerical revisions more directly. It applies the archived and latest risk values to the same REMX low-risk rule, the same March 2022 start, the same January 2022 frozen threshold, the same execution dates, the same market prices and the same costs. Only the risk observations used for classification differ.

The archived path ends at USD 13,844, a CAGR of 7.25\%, while the revised-data match ends at USD 18,550, or 14.22\%. Exactly two monthly investment classifications change. In both cases, the archived value puts the portfolio into REMX, while the revised value exceeds the low-risk threshold and leaves it in cash. Table~\ref{tab:attribution} identifies the decisions and the subsequent ETF returns.

\begin{table}[!htbp]
\centering\small
\caption{The two revised classifications and their wealth effects}
\label{tab:attribution}
\begin{tabular}{@{}lrrrr@{}}
\toprule
Decision & Risk: archive & Risk: revised & ETF return & Wealth effect \\
\midrule
2023-02-08 & 0.588 & 0.670 & -10.33 & 1,845 \\
2024-01-09 & 0.629 & 0.702 & -15.76 & 2,861 \\
\bottomrule
\end{tabular}
\tabnote{Decision dates are execution closes. Risk observations refer to the completed preceding month, and the common threshold is 0.651706 in original source units. ETF returns are gross holding returns until the next monthly execution close, in percent. Wealth effects are USD contributions to the matched terminal difference and include the portfolio's cash and cost conventions.}
\end{table}

The February 2023 decision uses a January risk value of about 0.588 in the archive and 0.670 in the revised file. REMX then loses about 10.33\% until the next execution close. The January 2024 decision uses a December value of about 0.629 versus 0.702, followed by a 15.76\% holding loss. Later revisions remove both losing exposures from the reconstructed low-risk strategy. The index values are source measurements, not probabilities of a particular geopolitical event.

Compounding makes a dollar attribution depend on the order in which the two decisions are changed. Let $W_{ab}$ denote final wealth when classification one is revised if $a=1$ and classification two if $b=1$. An equal average of changing decision one first and second gives
\begin{equation}
A_1=\tfrac12\bigl[(W_{10}-W_{00})+(W_{11}-W_{01})\bigr].
\label{eq:attribution}
\end{equation}
The corresponding average for decision two gives $A_2$, and the contributions add exactly to $W_{11}-W_{00}$. This allocates approximately USD 1,845 to the first decision and USD 2,861 to the second, totalling USD 4,706. It is a mechanical attribution with observed prices held fixed, not an estimate of an expected future revision profit.

\fig{09_vintage_audit}{Matched REMX wealth under archived and revised risk values. Shaded intervals mark the two holding periods with changed classifications. Sample, threshold, market prices, execution and costs are fixed; the difference comes from the risk version. The cash path provides the opportunity-cost comparison. Two avoided losing periods explain a large compounded difference despite most monthly classifications being unchanged.}{fig:vintages}

\section{Economic interpretation and limitations}

The physical-price and investment findings fit together. Geopolitical information can matter for a particular material quote during the month in which the news arrives. The controlled PrNd association and the local projections support studying that timing. A listed fund is a different claim, with operating costs, financing, country exposure and downstream demand added to the material-price channel. Acting after a monthly index is published also means investing after some price adjustment may already have occurred.

Economic value has several meanings here. Positive historical profit, improvement on passive ownership, improvement on a trend rule, lower forecast error and positive factor-adjusted alpha are different outcomes. The LIT and COPX avoidance rules deliver attractive revised-data historical CAGRs, but the uncertainty around their additional value is substantial. GPR-augmented forecasts can change profitable trading classifications without generally lowering squared errors. The archive exercise then shows that later revisions change those classifications again.

The availability distinction is especially important for the 2018 boundary. Sequential estimation with earlier closed return labels is an out-of-sample forecasting procedure with respect to coefficient estimation. It does not by itself guarantee that the data version existed then or that the strategy was chosen independently of the later results. The global working paper's publication also does not establish the first availability of every current country series. The complete 2018–2021 release sequence is not used here; the 2021 old-method snapshot is a sensitivity diagnostic, and the actual archived reconstruction begins in 2022.

The macro models remain conditional associations. They do not use an external geopolitical instrument, dated firm-level exposure or measured lost production to isolate a causal mechanism. Same-month controls may absorb part of a transmission channel, while lagged controls cannot remove every common influence. The IP comparison makes sample composition visible, and the response plot avoids claiming a precise persistent effect from wide intervals. The state definition and the original specification search remain material research choices.

Several practical boundaries also remain. Cash earns an interbank-rate proxy, trading costs are stated scenarios rather than observed historical spreads, and results exclude taxes, currency conversion, leverage and large-order market impact. US equity factors are informative but incomplete risk controls for globally exposed sector funds. The three-fund rotation benchmark has a deliberately narrow universe. None of these qualifications changes the reported arithmetic; they determine the economic scope of its interpretation.

For procurement, the PrNd result could motivate closer monitoring of inventories and suppliers during an already elevated-risk period. Actual savings would require a firm's dated purchase quotes, contracts, storage costs and delivery constraints. A research product could similarly combine material definitions, archived signals and decision-level attribution. Customer willingness to pay and commercial rights would need separate evidence; the paper does not convert investment backtests into a speculative sales valuation.

The most useful next validation is to fix a compact candidate family and evaluate genuinely new observations or an independent quote source. Additional production, inventory and export information could then test specific mechanisms. The present contribution is a reproducible account of where the existing evidence is strong, where profit depends on the implementation, and which decisions change when the information set changes.

\section{Conclusion}

The pricing and investment implications of geopolitical risk are state-dependent, commodity-specific and sensitive to historical information. In the selected PrNd specification, the high-state price association remains approximately three log-return percentage points after lagged economic controls. A matched industrial-production check gives similar controlled and uncontrolled estimates on its shorter sample. Local projections provide their clearest evidence for the current month, with substantially wider uncertainty about subsequent cumulative changes.

The chronological investment exercise begins in 2018, with initial estimation ending in 2017 and execution delayed for the older publication convention. Adding momentum, monthly and daily MA200, sector rotation, nested market-and-trend forecasts and six-factor attribution makes the investment comparison more demanding. Some geopolitical allocations earn substantial historical profits, but consistent incremental prediction and positive factor-adjusted alpha are not established after adjustment.

The archived 2022–2026 reconstruction adds a different contribution. Two revised monthly classifications remove losing REMX exposures and explain a USD 4,706 matched terminal-wealth difference. This demonstrates how a seemingly small change in a news index can materially alter a trading conclusion. Keeping product definition, timing, calibration, benchmark and data version explicit turns the evidence into testable research candidates rather than an automatic geopolitical trading rule.

\FloatBarrier
\begin{thebibliography}{99}
\small
\bibitem[Benjamini and Hochberg(1995)]{benjamini1995} Benjamini, Y., and Hochberg, Y. (1995). Controlling the false discovery rate: A practical and powerful approach to multiple testing. \textit{Journal of the Royal Statistical Society B}, 57(1), 289–300. \href{https://doi.org/10.1111/j.2517-6161.1995.tb02031.x}{doi:10.1111/j.2517-6161.1995.tb02031.x}.
\bibitem[Board of Governors(2026)]{fred2026} Board of Governors of the Federal Reserve System (2026). \textit{Federal Funds Effective Rate [DFF] and Nominal Broad U.S. Dollar Index [TWEXBGSMTH]}. FRED, Federal Reserve Bank of St. Louis. \href{https://fred.stlouisfed.org/series/DFF}{Cash-rate documentation}; \href{https://fred.stlouisfed.org/series/TWEXBGSMTH}{dollar-index documentation}.
\bibitem[Caldara and Iacoviello(2018)]{caldara2018} Caldara, D., and Iacoviello, M. (2018). Measuring geopolitical risk. \textit{International Finance Discussion Papers}, 1222, Federal Reserve Board, February. \href{https://doi.org/10.17016/IFDP.2018.1222}{doi:10.17016/IFDP.2018.1222}.
\bibitem[Caldara and Iacoviello(2021)]{gprold2021} Caldara, D., and Iacoviello, M. (2021). \textit{Old Geopolitical Risk Index: Webpage and legacy-method workbook}. Frozen 2021 history, retrieved November 17, 2026. \url{https://www.matteoiacoviello.com/gpr2019.htm}.
\bibitem[Caldara and Iacoviello(2022)]{caldara2022} Caldara, D., and Iacoviello, M. (2022). Measuring geopolitical risk. \textit{American Economic Review}, 112(4), 1194–1225. \href{https://doi.org/10.1257/aer.20191823}{doi:10.1257/aer.20191823}.
\bibitem[Caldara and Iacoviello(2026)]{gprdata} Caldara, D., and Iacoviello, M. (2026). \textit{Geopolitical Risk Index: Data, documentation and archived monthly releases}. Downloaded November 17, 2026. \url{https://www.matteoiacoviello.com/gpr.htm}.
\bibitem[Campbell and Thompson(2008)]{campbell2008} Campbell, J. Y., and Thompson, S. B. (2008). Predicting excess stock returns out of sample: Can anything beat the historical average? \textit{Review of Financial Studies}, 21(4), 1509–1531. \href{https://doi.org/10.1093/rfs/hhm055}{doi:10.1093/rfs/hhm055}.
\bibitem[Clark and West(2007)]{clark2007} Clark, T. E., and West, K. D. (2007). Approximately normal tests for equal predictive accuracy in nested models. \textit{Journal of Econometrics}, 138(1), 291–311. \href{https://doi.org/10.1016/j.jeconom.2006.05.023}{doi:10.1016/j.jeconom.2006.05.023}.
\bibitem[Depraiter and Goutte(2025)]{depraiter2025} Depraiter, L., and Goutte, S. (2025). Geopolitical risk and clean energy investments: Exploring the role of rare earths. \textit{International Review of Financial Analysis}, 100, 103965. \href{https://doi.org/10.1016/j.irfa.2025.103965}{doi:10.1016/j.irfa.2025.103965}.
\bibitem[Depraiter et~al.(2025)]{magnets2025} Depraiter, L., Goutte, S., and Porcher, T. (2025). Geopolitical risk and the global supply of rare earth permanent magnets: Insights from China's export trends. \textit{Energy Economics}, 146, 108496. \href{https://doi.org/10.1016/j.eneco.2025.108496}{doi:10.1016/j.eneco.2025.108496}.
\bibitem[European Parliament and Council(2024)]{eu2024} European Parliament and Council (2024). Regulation (EU) 2024/1252 establishing a framework for ensuring a secure and sustainable supply of critical raw materials. \url{https://eur-lex.europa.eu/eli/reg/2024/1252/oj}.
\bibitem[Fama and French(2015)]{fama2015} Fama, E. F., and French, K. R. (2015). A five-factor asset pricing model. \textit{Journal of Financial Economics}, 116(1), 1–22. \href{https://doi.org/10.1016/j.jfineco.2014.10.010}{doi:10.1016/j.jfineco.2014.10.010}.
\bibitem[Fan et~al.(2023)]{fan2023} Fan, J. H., Omura, A., and Roca, E. (2023). Geopolitics and rare earth metals. \textit{European Journal of Political Economy}, 78, 102356. \href{https://doi.org/10.1016/j.ejpoleco.2022.102356}{doi:10.1016/j.ejpoleco.2022.102356}.
\bibitem[Federal Reserve Bank of Dallas(2026)]{dallas2026} Federal Reserve Bank of Dallas (2026). \textit{Index of Global Real Economic Activity [IGREA]}. Data and methodology, retrieved November 17, 2026. \url{https://www.dallasfed.org/research/igrea}.
\bibitem[French(2026)]{french2026} French, K. R. (2026). \textit{Data Library: Daily five-factor and momentum returns}. September 2026 CRSP source files, retrieved November 17, 2026. \url{https://mba.tuck.dartmouth.edu/pages/faculty/ken.french/data_library.html}.
\bibitem[Global X(2026a)]{globalxcopx} Global X (2026a). \textit{Copper Miners ETF (COPX): Fund information and investment objective}. \url{https://www.globalxetfs.com/funds/copx/}.
\bibitem[Global X(2026b)]{globalxlit} Global X (2026b). \textit{Lithium \& Battery Tech ETF (LIT): Fund information and investment objective}. \url{https://www.globalxetfs.com/funds/lit/}.
\bibitem[GUK Magnetics(2026)]{guk2026} GUK Magnetics (2026). \textit{Raw Material and Commodity Price Tracker}. Neodymium and samarium workbook series. \url{https://gukmagnetics.co.uk/raw-material-commodity-price-tracker/}.
\bibitem[Holm(1979)]{holm1979} Holm, S. (1979). A simple sequentially rejective multiple test procedure. \textit{Scandinavian Journal of Statistics}, 6(2), 65–70. \href{https://doi.org/10.2307/4615733}{doi:10.2307/4615733}.
\bibitem[IEA(2025)]{iea2025} International Energy Agency (2025). \textit{Global Critical Minerals Outlook 2025}. IEA, Paris. \url{https://www.iea.org/reports/global-critical-minerals-outlook-2025}.
\bibitem[IMF(2019)]{imfmethod} International Monetary Fund (2019, with subsequent additions). \textit{IMF Primary Commodity Price Index: Technical documentation}. \href{https://www.imf.org/-/media/files/research/commodityprices/monthly/pcps-technical-documentation.pdf}{IMF technical documentation}.
\bibitem[IMF(2026)]{imfdata} International Monetary Fund (2026). \textit{Primary Commodity Prices}. Downloaded November 17, 2026; last retained observation August 2026. \url{https://www.imf.org/en/Research/commodity-prices}.
\bibitem[Jordà(2005)]{jorda2005} Jordà, Ò. (2005). Estimation and inference of impulse responses by local projections. \textit{American Economic Review}, 95(1), 161–182. \href{https://doi.org/10.1257/0002828053828518}{doi:10.1257/0002828053828518}.
\bibitem[Koenker and Bassett(1978)]{koenker1978} Koenker, R., and Bassett, G. (1978). Regression quantiles. \textit{Econometrica}, 46(1), 33–50. \href{https://doi.org/10.2307/1913643}{doi:10.2307/1913643}.
\bibitem[Künsch(1989)]{kunsch1989} Künsch, H. R. (1989). The jackknife and the bootstrap for general stationary observations. \textit{Annals of Statistics}, 17(3), 1217–1241. \href{https://doi.org/10.1214/aos/1176347265}{doi:10.1214/aos/1176347265}.
\bibitem[Li et~al.(2023)]{li2023} Li, Z.-Z., Meng, Q., Zhang, L., Lobont, O.-R., and Shen, Y. (2023). How do rare earth prices respond to economic and geopolitical factors? \textit{Resources Policy}, 85, 103853. \href{https://doi.org/10.1016/j.resourpol.2023.103853}{doi:10.1016/j.resourpol.2023.103853}.
\bibitem[Moskowitz et~al.(2012)]{moskowitz2012} Moskowitz, T. J., Ooi, Y. H., and Pedersen, L. H. (2012). Time series momentum. \textit{Journal of Financial Economics}, 104(2), 228–250. \href{https://doi.org/10.1016/j.jfineco.2011.11.003}{doi:10.1016/j.jfineco.2011.11.003}.
\bibitem[Newey and West(1987)]{newey1987} Newey, W. K., and West, K. D. (1987). A simple, positive semi-definite, heteroskedasticity and autocorrelation consistent covariance matrix. \textit{Econometrica}, 55(3), 703–708. \href{https://doi.org/10.2307/1913610}{doi:10.2307/1913610}.
\bibitem[OECD(2026)]{oecd2026} Organisation for Economic Co-operation and Development (2026). \textit{China composite leading indicator [CHNLOLITOAASTSAM] and industrial production [CHNPRINTO01IXPYM]}. Retrieved from FRED, November 17, 2026. \href{https://fred.stlouisfed.org/series/CHNLOLITOAASTSAM}{CLI documentation}; \href{https://fred.stlouisfed.org/series/CHNPRINTO01IXPYM}{industrial-production documentation}.
\bibitem[SDM Magnetics(2026)]{sdm2026} SDM Magnetics (2026). \textit{PrNd Mischmetal Price Trend}. Publisher quote history attributed to Asian Metal. \url{https://www.sdmmagnetic.com/prnd-mischmetal-price-trend/}.
\bibitem[VanEck(2026)]{vaneck2026} VanEck (2026). \textit{VanEck Rare Earth and Strategic Metals ETF (REMX): Overview and fact sheet}. \href{https://www.vaneck.com/us/en/investments/rare-earth-strategic-metals-etf-remx/overview/}{Official fund documentation}.
\bibitem[Yahoo Finance(2026)]{yahoo2026} Yahoo Finance (2026). \textit{Historical adjusted prices for REMX, LIT, COPX, SPY and additional equity controls}. Downloaded November 17, 2026. \url{https://finance.yahoo.com/quote/REMX/history/}.
\end{thebibliography}

\end{document}
